\begin{figure}[!htbp]
\centering
\resizebox{0.55\textwidth}{!}{%
\begin{tikzpicture}[scale=1.0]
\node[font=\small, anchor=west] at (-4.6, 3.9) {(a)};
\node[font=\scriptsize, anchor=west] at (-3.9, 3.9) {section $z_1{+}z_2{+}z_3=0$};
% The three walls {z_i = z_j} all contain the diagonal z_1=z_2=z_3, so in a transverse
% section they are three concurrent lines through the origin, 60 degrees apart, and they
% cut the section into exactly six sectors, one per ordering of three coordinates.
% one dash style per wall, so that the walls differ by more than color
\draw[dashed, thick, blue!70] (0, -4.2) -- (0, 4.2) node[above, font=\scriptsize, blue!80!black] {$z_1{=}z_2$};
\draw[dash dot, thick, green!50!black] (-3.637, -2.1) -- (3.637, 2.1) node[right, font=\scriptsize, green!50!black] {$z_2{=}z_3$};
\draw[dash dot dot, thick, red!70] (3.637, -2.1) -- (-3.637, 2.1) node[left, font=\scriptsize, red!80!black] {$z_1{=}z_3$};
\fill[black!45] (0,0) circle (1.6pt);
% Sector labels, counter-clockwise from the 0 degree direction
\node[font=\scriptsize, fill=white, inner sep=2pt, rounded corners=2pt] at (4.5, -0.55) {$z_1{>}z_3{>}z_2$};
\node[font=\scriptsize, fill=white, inner sep=2pt, rounded corners=2pt] at (1.45, 2.512) {$z_1{>}z_2{>}z_3$};
\node[font=\scriptsize, fill=white, inner sep=2pt, rounded corners=2pt] at (-1.45, 2.512) {$z_2{>}z_1{>}z_3$};
\node[font=\scriptsize, fill=white, inner sep=2pt, rounded corners=2pt] at (-2.9, 0) {$z_2{>}z_3{>}z_1$};
\node[font=\scriptsize, fill=white, inner sep=2pt, rounded corners=2pt] at (-1.45, -2.512) {$z_3{>}z_2{>}z_1$};
\node[font=\scriptsize, fill=white, inner sep=2pt, rounded corners=2pt] at (1.45, -2.512) {$z_3{>}z_1{>}z_2$};
% Data point and stability circle, drawn to scale. The centered point (1,-1,0) sits on the
% 0 degree ray at distance sqrt(2) from the origin; its nearest wall is at 1/sqrt(2), so with
% delta_min(z) = 1 the certified radius 1/2 is 0.707 of that distance.
\fill[black] (2.6, 0) circle (2pt);
\draw[thick, dotted, black!40] (2.6, 0) circle (0.919);
\draw[black!50] (2.6, 0) -- (1.681, 0);
\node[font=\scriptsize, black!75, anchor=north] at (2.14, -0.07) {$\delta_{\min}(z)/2$};
\node[font=\scriptsize, anchor=west] at (3.6, 1.1) {$z = (3, 1, 2)$};
\node[font=\scriptsize, anchor=west, black!75] at (3.6, 0.7) {$\argsort(z) = [2,3,1]$};
\end{tikzpicture}%
}\\[8pt]
\begin{tikzpicture}[>=Stealth, lbl/.style={font=\scriptsize}]
\node[font=\small, anchor=west] at (-0.6, 1.6) {(b)};
% one number line of e projected scores; shaded: the certified radius delta_min/2 around each score
\node[lbl, anchor=west] at (0.0, 1.6) {projected scores $z_1, \dots, z_e$ on one line; each shaded band has half-width $\delta_{\min}(z)/2$};
\foreach \x in {0.9, 3.0, 5.2, 6.2, 8.9, 11.9} {
    \fill[black!10] (\x-0.5, 0.48) rectangle (\x+0.5, 0.72);
}
\draw[thick] (0.0, 0.6) -- (13.2, 0.6);
\foreach \x in {0.9, 3.0, 5.2, 6.2, 8.9, 11.9} {
    \fill (\x, 0.6) circle (2.2pt);
}
% a perturbation inside the radius: every score stays in its band, the ranking is unchanged
\node[lbl, blue!70, anchor=west] at (0.0, 1.28) {every score moves less than $\delta_{\min}(z)/2$: the ranking is unchanged};
\foreach \x/\d in {0.9/0.35, 3.0/-0.3, 5.2/0.35, 6.2/-0.35, 8.9/0.3, 11.9/-0.35} {
    \draw[->, thick, blue!70] (\x, 0.98) -- (\x+\d, 0.98);
}
% the smallest gap
\draw[<->, thick, black!50] (5.2, 0.3) -- (6.2, 0.3);
\node[lbl, anchor=north] at (5.7, 0.26) {$\delta_{\min}(z)$};
% a perturbation beyond the radius: two neighbors cross, the ranking changes
\draw[->, thick, red!70] (5.2, -0.36) -- (5.9, -0.36);
\draw[->, thick, red!70] (6.2, -0.52) -- (5.5, -0.52);
\node[lbl, red!70, anchor=west, align=left] at (6.6, -0.44) {two neighbors move more than $\delta_{\min}(z)/2$ and cross:\\ the ranking changes};
\end{tikzpicture}
\caption{\textbf{Permutation cones and the stability radius.} (a)~Up to the hyperplanes $\{z_i=z_j\}$, whose points the tie rule assigns, the argsort encoding partitions the projected space $\R^e$ into $e!$ open convex cones (Weyl chambers). They are drawn here for $e=3$ in the plane $z_1+z_2+z_3=0$, which every cone meets and which the three hyperplanes cut into six sectors. All points in a cone map to the same permutation. The point $z=(3,1,2)$ is drawn at its position in the plane, which differs from $z$ by a multiple of the all-ones vector and has the same ranking and the same coordinate gaps. The dotted circle has radius $\delta_{\min}(z)/2$, half the smallest coordinate gap of Proposition~\ref{prop:score-stability}. It lies inside the region $\|\varepsilon\|_\infty<\delta_{\min}(z)/2$, so a perturbation inside it cannot cross a boundary. The nearest wall is farther away, at distance $\delta_{\min}(z)/\sqrt2$. (b)~The same condition on one line of projected scores. Scores that each move by less than $\delta_{\min}(z)/2$ keep their ranking, and two neighbors that move further can cross.}
\label{fig:cones}
\end{figure}
